%%%%%%%%%%%%%%%%%%%%%%%%%%%%%%%%%%%%%%%%%%%%%%%%%%%%%%%
% NOTAS TEORIA DE TRANSPORTE OPTIMO
% FCEN-UBA, segundo cuatrimestre de 2026. Fuente unica de las notas
% y de las guias de ejercicios (herramientas/generar_practicas.py).
%%%%%%%%%%%%%%%%%%%%%%%%%%%%%%%%%%%%%%%%%%%%%%%%%%%%%%%

\documentclass[11pt,a4paper]{amsart}
\usepackage{xr-hyper}
\newcounter{chapter}
\usepackage{fancyhdr}
\usepackage{geometry} \geometry{top=2cm,bottom=2.5cm,left=2.5cm,right=2.5cm}

\usepackage{amssymb}
\usepackage{hyperref}
\usepackage{listings}
\usepackage{mathtools}
\usepackage{enumitem}

% Fuzz -------------------------------------------------------------------
% Line spacing -----------------------------------------------------------

%%%%%%%%%fin pegada%%%%%%%%%%%%%%%%%%5

\usepackage[T1]{fontenc}
\usepackage[spanish]{babel}
\usepackage{graphicx}
\usepackage{array}
\usepackage{tikz}

%% Colores del Instituto de Cálculo y del Departamento de Matemática (portada y enlaces)
\definecolor{icnavy}{HTML}{013274}
\definecolor{iccyan}{HTML}{029BDF}
\definecolor{icmag}{HTML}{E2017B}
\definecolor{dmred}{HTML}{A8000C}
\hypersetup{colorlinks=true,linkcolor=icnavy,citecolor=icnavy,urlcolor=icnavy}
\usetikzlibrary{babel}

%%%%%%%%%%%%%%%%%%%%%%%%%%%%%%%%%%%%%%%%%%%%%%%%%%%%%%%%%%%%%%%%%%%%%%%%%%
%% ESTILO COMÚN DE LAS FIGURAS TikZ
%% Todas las figuras se declaran con [fig] y no llevan scale=: el tamaño se
%% controla con las coordenadas, para que el texto de las etiquetas salga
%% siempre del mismo tamaño. La paleta deriva de los colores institucionales,
%% de modo que las figuras quedan en la misma gama que la portada y los
%% hipervínculos.
%%%%%%%%%%%%%%%%%%%%%%%%%%%%%%%%%%%%%%%%%%%%%%%%%%%%%%%%%%%%%%%%%%%%%%%%%%
\usetikzlibrary{arrows.meta,positioning,calc,patterns}

\colorlet{figmain}{icnavy}        % objeto principal (gráficos de funciones)
\colorlet{figacc}{icmag}          % objeto destacado (recta soporte, soporte del plan)
\colorlet{figalt}{iccyan}         % objeto secundario (pseudoinversa, segundo panel)
\colorlet{figwarn}{dmred}         % contraejemplos, lo que no puede pasar
\colorlet{figgrey}{black!55}      % anotaciones, construcciones auxiliares
\colorlet{figbg}{icnavy!6}        % fondo de las regiones X x Y
\colorlet{figfill}{iccyan!15}     % relleno de regiones destacadas

\tikzset{
  fig/.style={
    scale=1,
    line cap=round, line join=round,
    every node/.append style={font=\footnotesize},
    >={Stealth[length=5pt,width=4pt]},
  },
  eje/.style      ={->, thin, black!75},
  rotulo/.style   ={font=\footnotesize},
  anota/.style    ={font=\scriptsize, figgrey},
  graf/.style     ={very thick, figmain},
  grafalt/.style  ={very thick, figalt},
  soporte/.style  ={thick, figacc},
  destaca/.style  ={line width=3pt, figacc, line cap=round},
  aux/.style      ={dashed, thin, figgrey},
  grilla/.style   ={very thin, black!18},
  region/.style   ={fill=figbg},
  regionA/.style  ={fill=figfill},
  regionB/.style  ={fill=figwarn!12},
  prohibido/.style={thick, dotted, figwarn},
  pt/.style       ={circle, fill, inner sep=0pt, minimum size=3.6pt},
  ptvacio/.style  ={circle, draw, thick, fill=white, inner sep=0pt, minimum size=3.6pt},
  ptmasa/.style   ={circle, fill=figacc, inner sep=0pt},
}
\newcommand{\figpanelsep}{6cm}   % separación única entre paneles de una figura



\parskip 3pt

%%
%% OPERADORES
%%
\renewcommand{\Re}{\operatorname{Re}}
\newcommand{\tr}{\operatorname{tr}}
\newcommand{\id}{\operatorname{id}}
\def\div{\operatorname {\mathrm{div}}}
\def\Dom{\operatorname {\mathrm{Dom}}}
\def\argmin{\operatorname {\mathrm{arg\, min}}}
\def\argmax{\operatorname {\mathrm{arg\, max}}}
\def\gen{\operatorname {\mathrm{gen}}}
\def\Im{\operatorname {\mathrm{Im}}}
\def\Nu{\operatorname {\mathrm{Nu}}}
\def\sop{\operatorname {\mathrm{sop}}}
\def\diam{\operatorname {\mathrm{diam}}}
\def\dist{\operatorname {\mathrm{dist}}}
\def\var{\operatorname {\mathrm{Var}}}
\def\codim{\operatorname {\mathrm{codim}}}
\def\pint{\operatorname {--\!\!\!\!\!\int\!\!\!\!\!--}}

%%
%% EPSILON
%%
\newcommand{\ve}{\varepsilon}

%%
%% LETRAS MATHBB
%%
\newcommand{\R}{{\mathbb R}}
\newcommand{\Z}{{\mathbb Z}}
\newcommand{\C}{{\mathbb C}}
\newcommand{\Q}{{\mathbb Q}}
\newcommand{\N}{{\mathbb N}}
\newcommand{\E}{{\mathbb E}}
\newcommand{\V}{{\mathbb V}}
\renewcommand{\P}{{\mathbb P}}
\newcommand{\II}{{\mathbb I}}
\newcommand{\OO}{{\mathbb O}}
\newcommand{\Sb}{{\mathbb S}}
\newcommand{\Tb}{{\mathbb T}}

%%
%% LETRAS CALIFGRAFICAS
%%
\newcommand{\F}{{\mathcal F}}
\newcommand{\Hc}{{\mathcal H}}
\newcommand{\Tc}{{\mathcal T}}
\newcommand{\Ec}{{\mathcal E}}
\newcommand{\Vc}{{\mathcal V}}
\newcommand{\A}{{\mathcal A}}
\newcommand{\J}{{\mathcal J}}
\newcommand{\G}{{\mathcal G}}
\newcommand{\U}{{\mathcal U}}
\newcommand{\B}{{\mathcal B}}
\newcommand{\Sw}{{\mathcal S}}
\newcommand{\D}{{\mathcal D}}
\newcommand{\M}{{\mathcal M}}
\newcommand{\LL}{{\mathcal L}}

%%
%% LETRAS MATHBF
%%
\newcommand{\ab}{{\mathbf a}}
\newcommand{\bb}{{\mathbf b}}
\newcommand{\cc}{{\mathbf c}}
\newcommand{\ee}{{\mathbf e}}
\newcommand{\ff}{{\mathbf f}}
\newcommand{\gb}{{\mathbf g}}
\newcommand{\hh}{{\mathbf h}}
\newcommand{\n}{{\mathbf n}}
\newcommand{\mm}{{\mathbf m}}
\newcommand{\pp}{{\mathbf p}}
\newcommand{\qq}{{\mathbf q}}
\newcommand{\sbf}{{\mathbf s}}
\newcommand{\tb}{{\mathbf t}}
\newcommand{\uu}{{\mathbf u}}
\newcommand{\vv}{{\mathbf v}}
\newcommand{\ww}{{\mathbf w}}
\newcommand{\xx}{{\mathbf x}}
\newcommand{\yy}{{\mathbf y}}
\newcommand{\zz}{{\mathbf z}}
\newcommand{\I}{{\mathbf 1}}
\newcommand{\Pb}{{\mathbf P}}
\newcommand{\XX}{{\mathbf X}}
\newcommand{\YY}{{\mathbf Y}}

%%
%% LETRAS MATHFRAK
%%

\newcommand{\Gk}{{\mathfrak G}}


%%
%% TEOREMAS
%%
\newtheorem{teo}{Teorema}[section]
\newtheorem{lema}[teo]{Lema}
\newtheorem{prop}[teo]{Proposición}
\newtheorem{cor}[teo]{Corolario}

%%
%% REMARK
%%
\theoremstyle{remark}
\newtheorem{obs}[teo]{Observación}

%%
%% DEFINICION
%%
\theoremstyle{definition}
\newtheorem{defi}[teo]{Definición}
\newtheorem{ejem}[teo]{Ejemplo}
\newtheorem{ejer}[teo]{Ejercicio}
\newtheorem{nota}[teo]{Notación}
\newtheorem{ejercicio}{Ejercicio}[chapter]


%%
%% NUMERACION
%%
\numberwithin{subsection}{section}
\numberwithin{equation}{section}

%%
%% TITULO
%%


\vfuzz7pt \hfuzz7pt
\externaldocument[N-]{../notas/Notas-TO}
\newcommand{\encabezado}[1]{%
\noindent \includegraphics[scale=2.1]{logo-dm-wide.png} \hskip 9.5cm
\includegraphics[scale=0.08]{logo-exactas-uba.jpg}
\smallskip
\usefont{T1}{phv}{m}{n}
\begin{center}{\Large\sc Teoría de Transporte Óptimo}\end{center}
\bigskip
\noindent{\sc #1}\\[-0.4cm]
\hrule
\normalfont}
\newcommand{\cuadernolink}[2]{\noindent\textbf{Cuaderno:} \emph{#1} (\href{https://colab.research.google.com/github/jfbonder/transporte-optimo/blob/main/cuadernos/#2}{abrir en Colab}).\par\medskip}

\begin{document}
\renewcommand{\thechapter}{11}
\encabezado{Práctica 11: La desigualdad de Brunn--Minkowski}
\bigskip
\pagestyle{fancy}
\fancyhead[L]{{\footnotesize{\usefont{T1}{phv}{m}{n}{\sc Teoría de Transporte Óptimo}}}}
\fancyhead[R]{\footnotesize \usefont{T1}{phv}{m}{n}{\sc Práctica 11}}
\thispagestyle{plain}
\noindent{\footnotesize Los ejercicios son los de la sección de ejercicios del Capítulo~11 de las notas del curso, con la misma numeración; las referencias a teoremas, proposiciones y secciones remiten a las notas (\url{https://github.com/jfbonder/transporte-optimo}). ($\bigstar$) ejercicio de mayor dificultad o que excede el programa.}
\medskip


\begin{ejercicio}[Casos elementales]\label{ejer.BM.elementales}
\leavevmode
\begin{enumerate}[label=(\alph*)]
\item Probar la desigualdad de Brunn--Minkowski directamente para $d=1$
(intervalos y, en general, compactos de $\R$) y para cajas
\[
A=\prod_{i=1}^d[a_i,b_i],\qquad B=\prod_{i=1}^d[c_i,d_i]
\]
en $\R^d$ (usar la desigualdad entre las medias aritmética y geométrica).
\item Probar que hay igualdad si $A$ y $B$ son convexos y homotéticos
($B=\lambda A+v$, $\lambda>0$), y mostrar que la hipótesis de convexidad no
se puede omitir: para $A=B=[0,1]\cup[2,3]\subset\R$, $|A+A|=6>4=|2A|$.
\item Sea $A$ un cubo y $B$ un segmento de longitud $L>0$ (en cualquier
dirección).
Calcular $|A+B|$ y comprobar que es estrictamente mayor que $|A|$ aunque
$|B|=0$.
\end{enumerate}
\end{ejercicio}

\begin{ejercicio}[La forma multiplicativa]\label{ejer.BM.multiplicativa}
Probar que la desigualdad de Brunn--Minkowski es equivalente a su forma
multiplicativa \emph{sin exponentes}:
\[
|(1-t)A+tB|\ge|A|^{1-t}|B|^{t}\qquad\text{para todo }t\in[0,1]
\]
(para $A$, $B$ compactos de medida positiva). \emph{Sugerencia:} para
deducir la forma aditiva de la multiplicativa, aplicarla a
$A'=|A|^{-1/d}A$, $B'=|B|^{-1/d}B$ y un $t$ bien elegido; para la
recíproca, usar la desigualdad entre medias. Observar que la forma
multiplicativa no depende de la dimensión, lo que es útil para demostrarla
por inducción (desigualdad de Prékopa--Leindler).
\end{ejercicio}

\begin{ejercicio}[Una identidad de la matriz de cofactores]\label{ejer.BM.cofactores}
Sea $u\in C^3(\Omega)$ y $M(x)=\operatorname{cof}(D^2u(x))$ la matriz de
cofactores de su hessiano, de modo que $M\,D^2u=(\det D^2u)\,I$.
\begin{enumerate}[label=(\alph*)]
\item Probar que $\sum_i\partial_iM_{ij}=0$ para todo $j$ (cada columna de
$M$ es un campo de divergencia nula). \emph{Sugerencia:} escribir
$M_{ij}$ como un determinante de orden $d-1$ y usar la simetría de las
derivadas terceras de $u$.
\item Deducir la \emph{forma de divergencia} de la ecuación de
Monge--Ampère,
\[
\det D^2u=\frac1d\sum_{ij}\partial_i\bigl(M_{ij}\,\partial_ju\bigr).
\]
\item Usar (b) para dar otra demostración del Lema~\ref{N-lema.MA} en la
forma integral: si $T=\nabla u$ es un difeomorfismo de $\Omega$ sobre
$\Omega'$, entonces $\int_\Omega\det D^2u\,dx=|\Omega'|$.
\end{enumerate}
\end{ejercicio}

\begin{ejercicio}[La desigualdad isoperimétrica]\label{ejer.BM.isoperimetrica}
Sea $A\subset\R^d$ compacto con $|A|>0$ y sea
$\operatorname{Per}(A)=\lim_{\ve\to0^+}\frac{|A+\ve\overline{B(0,1)}|-|A|}{\ve}$
su contenido de Minkowski, que suponemos existe.
\begin{enumerate}[label=(\alph*)]
\item Deducir del Teorema~\ref{N-teo.BM} la desigualdad isoperimétrica
\[
\operatorname{Per}(A)\ge d\,\omega_d^{1/d}|A|^{(d-1)/d},
\]
completando los detalles de la Observación~\ref{N-obs.isoperimetrica}.
\item Verificar que para bolas hay igualdad y que para un cubo de lado
$\ell$ el contenido de Minkowski es $2d\,\ell^{d-1}$; comparar ambos
miembros en ese caso.
\item Para $d=2$, deducir la forma clásica $L^2\ge4\pi\,\text{Área}$ para
dominios con borde $C^1$ (aceptando que el contenido de Minkowski coincide
con la longitud del borde).
\end{enumerate}
\end{ejercicio}

\begin{ejercicio}[Convexidad por desplazamiento de la entropía]\label{ejer.BM.entropia}
Sean $\mu_0=\rho_0\LL^d$ y $\mu_1=\rho_1\LL^d$ con densidades continuas y
positivas en abiertos acotados $\Omega_0$, $\Omega_1$, y supongamos que el
mapa de Brenier $T=\nabla u$ entre ellas satisface la hipótesis (R) de la
Sección~\ref{N-sec.BM.demo} ($u\in C^2$, $T$ difeomorfismo). Sea
$\mu_t=(T_t)_\#\mu_0$ la interpolación de McCann, con densidad $\rho_t$.
\begin{enumerate}[label=(\alph*)]
\item Probar la ecuación de Monge--Ampère
$\rho_t(T_t(x))\det DT_t(x)=\rho_0(x)$.
\item Probar que la \emph{entropía}
$\mathcal H(\mu_t)=\int\rho_t\log\rho_t\,dx$ satisface
\[
\mathcal H(\mu_t)=\int\rho_0(x)\log\rho_0(x)\,dx-\int\rho_0(x)\log\det\bigl((1-t)I+tD^2u(x)\bigr)\,dx,
\]
y que $t\mapsto\log\det((1-t)I+tM)$ es cóncava para $M$ simétrica definida
positiva. Concluir que $t\mapsto\mathcal H(\mu_t)$ es convexa.
\item Probar, con el Lema~\ref{N-lema.det.concavo}, que $t\mapsto\int\rho_t^m\,dx$
es convexa si $m>1$ y cóncava si $1-\frac1d\le m<1$ (de modo que
$-\int\rho_t^m$ es convexa, como en la
Observación~\ref{N-obs.desplazamiento.convexo}). \emph{Sugerencia:}
$\int\rho_t^m=\int\rho_0^m\,(\det DT_t)^{1-m}$, y $(\det DT_t)^{1/d}$ es
cóncava y positiva en $t$.
\end{enumerate}
\end{ejercicio}

\begin{ejercicio}[$\bigstar$ Una desigualdad de tipo Sobolev]\label{ejer.BM.sobolev}
Sea $\Omega\subset\R^d$ convexo y acotado, $f\in C^1(\overline\Omega)$,
$r\in[1,\infty]$ y $r'$ su exponente conjugado ($\frac1r+\frac1{r'}=1$).
Sean $\mu=\rho_0\LL^d$ y $\nu=\rho_1\LL^d$ en $\mathcal P(\Omega)$ con
$\|\rho_0\|_{L^r},\|\rho_1\|_{L^r}\le C$, y supongamos, como en el
Ejercicio~\ref{ejer.BM.entropia}, que el mapa de Brenier $T$ entre ellas
satisface la hipótesis (R). Probar que
\[
\Bigl|\int_\Omega f\,d(\mu-\nu)\Bigr|\le C^{1/2}\,\|\nabla f\|_{L^{2r'}(\Omega)}\,W_2(\mu,\nu).
\]
(Para $r=\infty$ el exponente es $2r'=2$; para $r=1$ es $2r'=\infty$ y la
desigualdad se reduce a $W_1\le W_2$.)
\emph{Esquema:} escribir $\int f\,d(\nu-\mu)=\int_0^1\!\!\int\nabla f(T_t(x))\cdot(T(x)-x)\,d\mu\,dt$
con $T_t=(1-t)\id+tT$ y $\mu_t=(T_t)_\#\mu=\rho_t\LL^d$ (que está concentrada en
$\Omega$ por convexidad); aplicar Cauchy--Schwarz en $L^2(\mu)$, que da el
factor $\bigl(\int|T-\id|^2d\mu\bigr)^{1/2}=W_2(\mu,\nu)$ y el factor
$\bigl(\int|\nabla f|^2\rho_t\bigr)^{1/2}$; acotar este último por Hölder con
exponentes $(r',r)$, y controlar $\|\rho_t\|_{L^r}\le C$ uniformemente en $t$:
para $r<\infty$ por la convexidad de $t\mapsto\int\rho_t^r$
(Ejercicio~\ref{ejer.BM.entropia}(c)), y para $r=\infty$ por la ecuación de
Monge--Ampère y la concavidad de $t\mapsto\det DT_t^{1/d}$
(Lema~\ref{N-lema.det.concavo}). Más generalmente, para $q\in(1,\infty)$ y
$q'=q/(q-1)$ vale la misma desigualdad con $W_q$, $\|\nabla f\|_{L^{q'r'}}$
y la constante $C^{1/q'}$, usando el mapa óptimo para el costo $|x-y|^q$;
ver \cite[Sec.~5.5]{Santambrogio}.
\end{ejercicio}

\bibliographystyle{plain}
\bibliography{../notas/biblio}

\end{document}
